\section{Objetivos}

\subsection{Objetivo general}

Diseñar, implementar y evaluar una arquitectura distribuida para la coordinación de coaliciones multi-AMR en el transporte de cargas heterogéneas, integrando dinámicas poblacionales, cierre entero de quórum, capacidad efectiva en espacio \emph{wrench} y validación experimental reproducible por simulación.

\subsection{Objetivos específicos}

\begin{enumerate}[label=O\arabic*.,leftmargin=2.2em]
  \item Formular el transporte cooperativo como un problema de asignación física, donde cada robot decide localmente su participación en una carga y la utilidad depende de déficit, coste operativo y factibilidad mecánica.
  \item Construir una escalera de validación que avance desde reclutamiento escalar hasta capacidad heterogénea, contactos, movimiento, transporte, recuperación, comunicación y escala de flota.
  \item Diseñar Smith-QR como extensión de la dinámica de Smith con memoria local, estimación de ocupación, cierre entero de coaliciones y guardias de ejecución.
  \item Extender el criterio de quórum desde cardinalidad hacia capacidad efectiva, incorporando fuerza, torque, orientación, batería, distancia y conectividad.
  \item Comparar la arquitectura propuesta frente a referencias centralizadas, heurísticas clásicas, subastas distribuidas, variantes basadas en modelo y comparadores aprendidos.
  \item Garantizar trazabilidad experimental mediante semillas disjuntas, mundos compartidos, métricas homogéneas, contrastes pareados y límites de validez explícitos.
\end{enumerate}

\subsection{Alcance}

El alcance defendible de la memoria es simulación reproducible con integración visual y análisis cuantitativo. No se declara validación física en hardware, certificación de seguridad, percepción real ni despliegue industrial. Los métodos aprendidos se usan como comparadores y no como sustituto de la contribución analítica principal. La memoria distingue tres niveles de madurez: resultados implementados y evaluados en campañas principales, pruebas de consistencia matemática o sintética, y extensiones que quedan como trabajo futuro.

En particular, la asignación explícita de roles de contacto se aproxima mediante formación derivada dentro de cada coalición, no como una optimización completa de todos los pares robot--carga--slot. La batería se incorpora como restricción económica y de retornabilidad, sin modelar todavía una red completa de estaciones, colas y recarga física.